\section{Discussion}
\label{sec:discussion}

\subsection{Principal Findings}
\textbf{Escalation is necessary, and it can be made provably correct---but it cannot repair detection.} When the owner cannot respond, the ladder converts \R{sim.ownerincapacitated.p300givendet}\% of \emph{detected} alerts into a human response within five minutes in simulation, and its safety properties hold in every interleaving the model checker explored. Yet the end-to-end figure (\R{sim.ownerincapacitated.p300}\%) is set almost entirely by the detector's episode sensitivity (\R{sim.pdetect}). Everything downstream of the gate multiplies by that number. The cheapest available improvement is behavioural rather than computational: if the dog repeats an unacknowledged alert, detection rises to \R{sim.rep3.pdet}\% without any change to the model.

\textbf{Measured honestly, the false-alarm rate is far above what any escalation chain tolerates.} The owner-acknowledgement stage does real work: it stops roughly three in four false alarms (\R{burden.preach}\% get through with the shipped $T_1$). The detector, however, produces \R{cont.rfe.r60.fah}~FA/h on continuous training-session motion even with the refractory period, and \R{cont.rfe.r0.fah}~FA/h without it. Our burden analysis turns this into a requirement: \R{req.maxfah}~FA/h for one caregiver page per week, three orders of magnitude below what we measure, and comparable to what clinically validated seizure wearables achieve~\cite{naganur2022meta,meritamlarsen2023nonEEG}. That number, rather than a further increment in F1, is what we propose as the target for detection work in this application. The negatives in this dataset are the hardest available---a dog being actively cued, seconds from spinning on command---so the all-day rate is very probably lower. It is unmeasured, and the source study's 412.6\,h of unlabelled wear data~\cite{brady2025collar} is what would measure it.

\textbf{Curated-clip evaluation understates false alarms.} The protocol inherited from the source study scores fixed clips that keep away from spins, and implies \R{op.single.fah}~FA/h; scoring the same detector on continuous recordings gives \R{cont.ratio.curated}$\times$ that. A large part of the difference survives even generous post-spin tolerances. For any system whose output triggers an action, we recommend that this field adopt what seizure detection already standardised~\cite{dan2025szcore,beniczky2018standards}: continuous-recording, event-based scoring, false alarms normalised by exposure, and operating points chosen out of sample.

\textbf{On this data, learned models add calibration, not discrimination.} A one-statistic heuristic matches the forest's AUC, matches it on continuous streams once a refractory period is applied, and neither a CNN nor ROCKET helps. With 135 events from six dogs this is unsurprising, and it argues for collecting data before engineering models.

\textbf{Formal methods paid for themselves.} A bounded model of a small protocol found a design flaw (the watchdog bypassing the caregiver tier when the phone fails) that the existing unit tests and the end-to-end replay never exercised, and writing the model exposed an integration bug that disabled the watchdog in deployment while every test passed. Both are the kind of defect that surfaces only in an incident review; here they surfaced at design time. The heartbeat fix and its regression test are now part of the system; the Ladder watchdog is specified, model-checked and simulated, and porting it to the production backend is listed under future work. We suggest that any alerting system that splits its state across a device and a server is worth the same treatment.

\subsection{Relation to the Source Study}
Our LOGO segment F1 for the 200-tree forest (\R{logo.rf.segF1}) is close to the published LODO F1 of 0.65~\cite{brady2025collar}, and we reproduce the qualitative findings: supervised models beat the heuristic, within-subject results are optimistic, and event-level aggregation recovers most missed windows. The numbers are not directly comparable---our groups are not dogs (Section~\ref{sec:data}), and we fit selection inside every fold---and we do not claim that the differences in either direction are meaningful given the interval widths in Table~\ref{tab:logo}. Two dataset properties that we report---the 51.2\,Hz rate and the gyroscope saturation---should be accounted for by anyone reusing the data.

\subsection{Design Implications}
\begin{itemize}
\item \textbf{Make the detector rotation-robust before deployment}, by rotation augmentation or by restricting it to magnitude features; collar rotation was the only simulated transfer effect that mattered, and it mattered a great deal (\R{rob.c.rotx90deg.evRec} event recall after a 90\textdegree{} rotation).
\item \textbf{Configure the IMU range deliberately.} The MPU6050's default $\pm$250\,\textdegree{}/s range clips spins and costs recall; a range at least as wide as the training sensor's must be set, and clipping the wider range to the training range in firmware keeps the feature distribution matched.
\item \textbf{Replace the Direct watchdog with the ladder-aware one, and make paging idempotent} at a shared provider keyed by alert and channel, so that relay and backend cannot both page the same person.
\item \textbf{Give the collar a second path.} Silent loss---the phone dies before the backend hears of the alert---cannot be fixed on the server side; a local audible alarm on the collar, or a low-power cellular link, is the only remedy.
\item \textbf{Set $T_1$ per owner}, with their clinician: shortening it from 60\,s to 30\,s improves time-to-help for an incapacitated owner but lets \R{grid.t1short.preach}\% of false alarms through to the caregiver instead of \R{burden.preach}\%.
\end{itemize}

\subsection{Limitations}
\textbf{Data.} Six young dogs of three breeds from one training organisation, performing a cued behaviour in training sessions; no dog identities, so no true LODO; \R{ds.posevents} positive events, which makes every interval wide; negatives drawn only from training-session motion; \R{ds.hours}\,h of continuous recordings in total. Spontaneous alerts during real pre-ictal states may look different, and whether the dog's alert predicts a seizure at all is outside the scope of this work and of the data.

\textbf{Hardware.} The firmware has not run on an ESP32 and the collar has not been worn. The C core is verified against the Python pipeline on the host; timing, power, BLE behaviour and the real Shimmer3$\rightarrow$MPU6050 transfer are unmeasured, and the transfer analysis is a simulation built from datasheet tolerances, not a substitute for a paired recording.

\textbf{Simulation.} Every human-behaviour parameter is an assumption. The conclusions we draw from the simulation are therefore deliberately structural---outcomes are bounded by detection, the watchdog variant removes tier-skipping, the false-alarm requirement scales linearly with the rate---rather than numerical predictions of what would happen to a particular owner.

\textbf{Formal model.} The model covers one alert episode with discrete time and a bounded horizon. It does not model concurrent alerts, a relay process that restarts and loses its in-memory state, clock drift between relay and backend, or the notification providers themselves. The Java backend that implements the watchdog in production was syntax-checked but not built in this work; the watchdog behaviour simulated and model-checked here is a port of its logic.

\subsection{Future Work}
In order of expected value: (1) obtain the recording-to-dog mapping and re-run under true LODO; (2) measure false alarms on the unlabelled all-day wear data; (3) record a dog performing cued spins with a Shimmer3 and an MPU6050 collar simultaneously; (4) establish whether trained dogs repeat an unacknowledged alert, and train for it if not; (5) implement the ladder-aware watchdog and idempotent paging in the production backend and extend the formal model to concurrent alerts and relay restarts; and only then (6) a supervised trial with a handler and a clinician, in which the system is explicitly not relied upon.
